\section*{अध्याय \ref{ch:g10:sport-and-health} --- शारीरिक क्रिया और स्वास्थ्य}

\begin{solution}{exo:g10:sport-and-health:1}
बड़ा हृदय, बड़े प्रघात आयतन और कम विश्राम दर के साथ; \omterm{def:g10:muscles-and-joints:muscle}{पेशी तंतुओं} में अधिक
केशिकाएँ और \omterm{def:g10:cells-common-unit:organelle}{सूत्रकणिकाएँ}; अधिक लाल रक्त \omterm{def:g10:cells-common-unit:cell}{कोशिकाएँ}; बड़े \omterm{def:g10:exercise-and-energy:glycogen}{ग्लाइकोजन} भंडार
(और रक्त-शर्करा तथा रक्तचाप का बेहतर नियंत्रण)।
\end{solution}

\begin{solution}{exo:g10:sport-and-health:2}
प्रघात आयतन बढ़ चुका है, इसलिए वही \qty{5}{L/min} कम धड़कनों में पहुँच
जाता है: $Q = f \times V_s$ में बड़े $V_s$ के लिए छोटा $f$ ही काफ़ी
है।
\end{solution}

\begin{solution}{exo:g10:sport-and-health:3}
दिन में लगभग एक घंटा मध्यम क्रिया। नहीं: चलना, स्कूल तक साइकिल, सीढ़ियाँ
और घर का काम, सब गिने जाते हैं --- मायने वह हलचल रखती है जो \omterm{def:g10:exercise-and-energy:expenditure}{ऊर्जा व्यय}
बढ़ा दे।
\end{solution}

\begin{solution}{exo:g10:sport-and-health:4}
उन पदार्थों या तरीक़ों का इस्तेमाल जो अनुकूलन बनाने के बजाय शरीर के
नियंत्रण को दरकिनार करके प्रदर्शन बढ़ाते हैं। EPO लाल \omterm{def:g10:cells-common-unit:cell}{कोशिकाओं} की संख्या
के नियंत्रण को दरकिनार करता है; और उत्तेजक थकान को, यानी उस चेतावनी को जो
हृदय की रक्षा करती है।
\end{solution}

\begin{solution}{exo:g10:sport-and-health:5}
टिकी हुई थकान, चढ़ती हुई विश्राम की हृदय दर, गिरता प्रदर्शन (और साथ ही
बिगड़ी नींद, बार-बार संक्रमण, \omterm{def:g10:muscles-and-joints:muscle}{कंडरा} या हड्डी में दर्द)।
\end{solution}

\begin{solution}{exo:g10:sport-and-health:6}
12 हफ़्तों के बाद: चार सत्रों के लिए $72 - 60 = 12$ धड़कनें, और दो के लिए
$72 - 64 = 8$। अनुकूलन माँग के अनुपात में होता है --- अधिक सत्र, अधिक
बदलाव --- और दोनों समूहों में वह नई अवस्था के पास पहुँचते-पहुँचते धीमा पड़
जाता है।
\end{solution}

\begin{solution}{exo:g10:sport-and-health:7}
$V_s = Q/f$: $5000/72 \approx \qty{69}{mL}$ से $5000/54 \approx
\qty{93}{mL}$, यानी $72/54 \approx 1.33$ का गुणक।
\end{solution}

\begin{solution}{exo:g10:sport-and-health:8}
$1 - 0.55 = 45\%$ कम। सबसे बड़ा लाभ पहला ही क़दम देता है, सबसे कम सक्रिय
से कम तक (1.00 से 0.80)।
\end{solution}

\begin{solution}{exo:g10:sport-and-health:9}
\qty{2.0}{kg}, यानी उसके द्रव्यमान का $2.0/55 \approx 3.6\%$। उसका
रक्त-आयतन गिरता है, निर्गम बनाए रखने के लिए हृदय दर चढ़ती है, शरीर का ताप
बढ़ता है, प्रदर्शन गिरता है; और ऐंठन की संभावना बन जाती है।
\end{solution}

\begin{solution}{exo:g10:sport-and-health:10}
$10 \times 1.1^7 \approx \qty{19.5}{km}$, इसलिए \qty{20}{km} आठवें
हफ़्ते में पहुँचता है: यानी लगभग दो महीने।
\end{solution}

\begin{solution}{exo:g10:sport-and-health:11}
हर तंतु पेशीतंतुक जोड़कर मोटा होता है, और तंत्रिका तंत्र एक साथ तथा बेहतर
तालमेल से अधिक तंतु भर्ती करना सीख लेता है; दोनों बल बढ़ाते हैं। तंतुओं की
संख्या वयस्कावस्था में तय रहती है।
\end{solution}

\begin{solution}{exo:g10:sport-and-health:12}
प्रति लीटर ऑक्सीजन $56/45 \approx 1.24$ गुना, यानी लगभग 24\% बढ़ती है;
वही \omterm{def:g10:heart-lungs-effort:output}{हृदय-निर्गम} और वही निष्कर्षण मानें तो अधिकतम $\dot V\!\mathrm{O_2}$
भी उतने ही अंश तक बढ़ सकती है। पर 56\% \omterm{def:g10:cells-common-unit:cell}{कोशिकाओं} वाला रक्त कहीं गाढ़ा है;
और सोते हृदय से धीमा पड़ा वह हृदय या मस्तिष्क की किसी वाहिका में जम सकता
है।
\end{solution}

\begin{solution}{exo:g10:sport-and-health:13}
ऊँचाई पर शरीर ख़ुद ऑक्सीजन की असली कमी के जवाब में, एक नपी-तुली मात्रा
में, अपना EPO बढ़ाता है, और कमी ख़त्म होते ही उसे दोबारा घटा देता है:
यानी नियंत्रण काम पर है। इंजेक्ट किया गया EPO वह संख्या थोप देता है जो
नियंत्रण कभी चुनता ही नहीं, न उसे उचित ठहराने वाली कोई कमी होती है, न कोई
ब्रेक।
\end{solution}

\begin{solution}{exo:g10:sport-and-health:14}
काम करती पेशी ग्लूकोज़ बिना इंसुलिन के ले लेती है, और प्रशिक्षित पेशी उसके
प्रति अधिक संवेदनशील बनी रहती है; इसलिए चलना रक्त-शर्करा सीधे घटाता है और
हर भोजन पर ज़रूरी इंसुलिन कम कर देता है। दवा शर्करा तो घटा देती है, पर
कारण --- यानी असंवेदनशील, बिना इस्तेमाल की पेशी --- वहीं का वहीं छोड़
देती है, और हृदय, हड्डियों तथा द्रव्यमान के लिए कुछ नहीं करती, जिन्हें
चलना सुधार देता है।
\end{solution}

\begin{solution}{exo:g10:sport-and-health:15}
अनुकूलन दोहराई गई, धीरे-धीरे बढ़ाई गई माँगों के बीच के विश्राम में बनता
है: सघन भार का एक महीना चुस्ती के बजाय अति-प्रशिक्षण और चोट देने की अधिक
संभावना रखता है, क्योंकि \omterm{def:g10:muscles-and-joints:muscle}{कंडरा} और हड्डी हफ़्तों में नहीं, महीनों में ढलते
हैं। और रुकते ही ये कमाई हफ़्तों में खुल जाती है: किसी एक झोंके से साल भर
के लिए कुछ भी ``तैयार'' नहीं होता। साल भर की नियमित मध्यम क्रिया वह कर
देती है जो वह महीना नहीं कर सकता।
\end{solution}

\begin{solution}{pb:g10:sport-and-health:1}
\textbf{1.} पहले $4800/72 \approx \qty{67}{mL}$; बाद में $4800/56 \approx
\qty{86}{mL}$।

\textbf{2.} $42 \times 62 \approx \qty{2.6}{L/min}$; $50 \times 62 =
\qty{3.1}{L/min}$: यानी 19\% का लाभ।

\textbf{3.} अधिकतम निर्गम $\dot V\!\mathrm{O_2}/0.150$: \num{17.4} से
\qty{20.7}{L/min}; और 200 धड़कनों पर अधिकतम प्रघात आयतन \num{87} से
\qty{103}{mL}।

\textbf{4.} बड़ा, मज़बूत हृदय प्रघात आयतन को समझाता है; और तंतुओं में
अधिक केशिकाएँ तथा \omterm{def:g10:cells-common-unit:organelle}{सूत्रकणिकाएँ} निष्कर्षण के लाभ को समझातीं।

\textbf{5.} 26 हफ़्तों में 4 का गुणक औसतन $4^{1/26}
\approx 1.055$ है, यानी हफ़्ते में लगभग 5.5\% --- जो 10\% के नियम के भीतर
है ($1.1^{26}
\approx 12$ इससे कहीं अधिक की छूट देता)।

\textbf{6.} अति-प्रशिक्षण: चढ़ती विश्राम की हृदय दर, बिगड़ी नींद,
संक्रमण, गिरता प्रदर्शन, और धीरे मरम्मत होने वाले किसी ऊतक में दर्द।

\textbf{7.} पिंडली की हड्डी की तनाव-चोट (कोई तनाव-भंग या उसकी सतह की
सूजन): हड्डी सब ऊतकों में सबसे धीरे ढलती है, और भार पखवाड़े भर में दोगुना
कर दिया गया था।

\textbf{8.} एक सत्र से अगले तक उबरना अधूरा रह जाता है; शरीर टिके हुए
तनाव की अवस्था में है, और जिस अनुकूलन ने दर घटाई थी वह उलटा जा रहा है।

\textbf{9.} लगभग विश्राम और नींद का एक हफ़्ता (अनुकूलन विश्राम में बनता
है); फिर पिछली दूरी की आधी पर दो हफ़्ते, कोई रफ़्तार का काम नहीं, और
हफ़्ते में दसवें हिस्से भर की दोबारा बढ़त; पिंडली किसी डॉक्टर को दिखाई
जाए और दर्द ख़त्म होने तक आराम दिया जाए (दर्द सुनिए); और मरम्मत के लिए
पीना तथा खाना।

\textbf{10.} एक उत्तेजक: वह थकान को ढँक देता है, यानी उस चेतावनी को जो
हृदय को निढाल होने और अति-गर्मी से बचाती है।

\textbf{11.} सैम $32 \times 80 \approx \qty{2.6}{L/min}$; एलेक्स $48
\times 64 \approx \qty{3.1}{L/min}$। एलेक्स की निरपेक्ष रूप में भी बड़ी
है; और रोज़मर्रा का जीवन --- सीढ़ियाँ, चलना, ख़ुद को ढोना --- शरीर का
अपना ही द्रव्यमान हिलाता है, इसलिए प्रति किलोग्राम वाला आँकड़ा ही तय करता
है कि कैसा महसूस होगा।

\textbf{12.} सैम अपनी अधिकतम का $12/32 \approx 38\%$, और एलेक्स $12/48 =
25\%$। सैम, जो अपनी छत के एक-तिहाई से ऊपर काम कर रहा है, हाँफता हुआ
पहुँचता है।

\textbf{13.} बिना इस्तेमाल की पेशी रक्त से कम ग्लूकोज़ लेती है और
इंसुलिन के प्रति अपनी संवेदनशीलता खो देती है; उन्हीं भोजनों के लिए हर साल
अधिक इंसुलिन चाहिए और रक्त-शर्करा चढ़ने लगती है।

\textbf{14.} 1.00 से 0.70 तक: यानी लगभग 30\% कम ख़तरा।

\textbf{15.} \omterm{def:g10:sport-and-health:training}{शारीरिक क्रिया} वह हर हलचल है जो \omterm{def:g10:exercise-and-energy:expenditure}{ऊर्जा व्यय} बढ़ाए: काम तक
चलकर जाना, लिफ़्ट के बजाय सीढ़ियाँ, घर का काम --- दिन में उसका आधा घंटा न
किसी खेल की माँग करता है, न किसी साज़-सामान की।

\textbf{16.} हर साल 1\% की दर से 10\% का हाशिया दस साल की गिरावट है।

\textbf{17.} सैम की चोटी नीची थी और वह पच्चीस साल से हड्डी खो रहा है;
उसकी कलाई उस दहलीज़ के पार जा चुकी है जिस पर गिरने से वह टूट जाती है,
जबकि एलेक्स की हड्डी, जो किशोरावस्था में भार के नीचे बनी और तब से बनी रही,
दस साल आगे है।

\textbf{18.} अनुकूलन इस्तेमाल से बनी रहने वाली कोई रचना है, कोई स्थायी
कमाई नहीं: माँग रुकते ही वह खोल दी जाती है और माँग लौटते ही दोबारा बना दी
जाती है।

\textbf{19.} हृदय और धमनी के रोग, टाइप 2 मधुमेह, अस्थिसुषिरता (और साथ ही
कुछ कैंसर तथा अवसाद)। मधुमेह: निष्क्रिय पेशी इंसुलिन के प्रति अपनी
संवेदनशीलता खो देती है, रक्त-शर्करा चढ़ती है और अग्न्याशय पर बोझ पड़ता
है।

\textbf{20.} विश्राम की हृदय दर 78 के मुक़ाबले 54; अधिकतम
$\dot V\!\mathrm{O_2}$ प्रति किलोग्राम \qty{32}{mL/min} के मुक़ाबले
\num{48}; और चौदह साल तक निभाई गई हफ़्ते में लगभग तीन घंटे की दौड़।
\end{solution}
